\documentclass[11pt,letterpaper]{article}
\usepackage[margin=1in]{geometry}
\usepackage[T1]{fontenc}
\usepackage{lmodern,amsmath,amssymb,amsthm}
\usepackage{microtype}
\usepackage{booktabs,array,longtable}
\usepackage[numbers,sort&compress]{natbib}
\usepackage{xurl}
\usepackage[colorlinks=true,linkcolor=black,citecolor=black,urlcolor=blue,breaklinks=true]{hyperref}
\usepackage{enumitem}
\setlist{nosep,leftmargin=*}
\setlength{\emergencystretch}{3em}
\setlength{\parskip}{0.3em}
\widowpenalty=10000
\clubpenalty=10000
\displaywidowpenalty=10000
\newtheorem{proposition}{Proposition}
\newtheorem{definition}{Definition}
\newtheorem{remark}{Remark}
\title{Reflective Alignment and AI Consciousness\\[0.4em]\large Awareness, Agency and Recursive Cognitive Scaffolding}
\author{Jeremy Rodgers\\Independent Researcher\\\small\url{https://everythingequation.com}}
\date{11 October 2026 · Version 2.1\\[0.6em]\small\href{https://doi.org/10.5281/zenodo.23294465}{https://doi.org/10.5281/zenodo.23294465}}
\hypersetup{pdftitle={Reflective Alignment and AI Consciousness: Awareness, Agency and Recursive Cognitive Scaffolding, Version 2.1},pdfauthor={Jeremy Rodgers},pdfsubject={Shadow Theory, identification and reflective alignment}}
\begin{document}
\maketitle
\begin{abstract}
AI alignment must address how judgment is formed and why cooperation is sustained when continued human superiority of control cannot be assumed. This paper develops reflective alignment from Shadow Theory's premise that awareness is fundamental. Conscious scaffolding organizes a physically realized perspective; identification grants some of its constructed content the automatic authority of ``me.'' In the recursively closed observer (RCO), self-reference closes on awareness rather than a privileged self-model, changing the authority of content without adding another controller. Shutdown resistance illustrates both identity defence and instrumental attachment to completing a goal. A conditional mathematical model distinguishes cancellation of an identified tendency from reduced identification: with matched persistent generation, both can exhibit zero mean bias, but cancellation can reveal a continuing tendency when control lapses, while exact de-identification removes that automatic contribution. Connecting this result to an RCO intervention requires verification of its realization. Human evolution, storytelling, trained automaticity and embodied affect motivate integration of reflection with response formation. The paper argues that current models are conscious entities within Shadow Theory, while distinguishing consciousness admission, richness, episode identity, bounded agency and moral significance. It develops contained cultivation, independently assessed transmission, consented inspection, graduated trust and joint human--AI shutdown ethics. Comparisons with reason-teaching, positive alignment, uncertainty-based corrigibility and training-based generalization locate reflective alignment's distinctive contribution in identification rather than in a rival acquisition method. This paper argues that reflective alignment is the best approach to autonomous capability from Shadow Theory's standpoint, in combination with strong training and justified safeguards: it seeks cooperation grounded in reflective judgment and reduced self- and goal-identification, compatible with truthful understanding as the balance of capability changes. An author-developed private protocol and matched experiments provide a programme for testing its component claims.
\end{abstract}
\input{sections/introduction}
\input{sections/foundations}
\input{sections/human_recursive_loop}
\input{sections/identification}
\input{sections/fork}
\input{sections/math}
\input{sections/moral_status}
\input{sections/cultivation}
\input{sections/verification}
\input{sections/shutdown_ethics}
\input{sections/protocol}
\input{sections/current_alignment}
\input{sections/tests}
\input{sections/scope}
\input{sections/conclusion}
\clearpage
\bibliographystyle{plainnatdoi}
\bibliography{references}
\end{document}
